\documentclass[10pt,leqno]{amsart}
\usepackage{graphicx}
\baselineskip=16pt

\usepackage{indentfirst,csquotes}

\topmargin= .5cm
\textheight= 22.5cm
\textwidth= 32cc
\baselineskip=16pt

\evensidemargin= .9cm
\oddsidemargin= .9cm

\usepackage{amssymb,amsthm,amsmath}
\usepackage{xcolor,paralist,hyperref,titlesec,fancyhdr,etoolbox}
\newtheorem{theorem}{Theorem}[]
\newtheorem{definition}[theorem]{Definition}
\newtheorem{example}[theorem]{Example}
\newtheorem{lemma}[theorem]{Lemma}
\newtheorem{proposition}[theorem]{Proposition}
\newtheorem{corollary}[theorem]{Corollary}
\newtheorem{conjecture}[theorem]{Conjecture}

\titleformat{\section}[display]{\normalfont\huge\bfseries\centering}{\centering\chaptertitlename\thechapter}{10pt}{\Large}
\titlespacing*{\section}{0pt}{0ex}{0ex}

\hypersetup{ colorlinks=true, linkcolor=black, filecolor=black, urlcolor=black }
\def\proof{\noindent {\it Proof. $\, $}}
\def\endproof{\hfill $\Box$ \vskip 5 pt }

% --- Additions (same as report.tex) -----------------------------------------
\usepackage{booktabs}
\newcommand{\heading}[1]{\medskip\noindent\textbf{#1}\par}
\newcommand{\qa}[2]{\smallskip\noindent\textbf{Q.} \emph{#1}\par\noindent\textbf{A.} #2\par}
% Generated by src/report_assets.py -- do not edit by hand.
\newcommand{\DataNiftyDays}{4{,}740}
\newcommand{\DataVixMonths}{196}
\newcommand{\DataVixVerified}{4{,}032}
\newcommand{\DataYahooMissing}{47}
\newcommand{\FoAtmViolPct}{0.28}
\newcommand{\FoDays}{4{,}031}
\newcommand{\FoExpiries}{194}
\newcommand{\FoIVcorr}{0.96}
\newcommand{\FoParityBand}{28}
\newcommand{\FoParityPairs}{136{,}118}
\newcommand{\IvrOptAbove}{73}
\newcommand{\IvrOptIvMean}{15.7}
\newcommand{\IvrOptRvMean}{14.5}
\newcommand{\IvrOptVarRatioCovid}{0.68}
\newcommand{\IvrOptVarRatioEx}{1.20}
\newcommand{\IvrVixAbove}{79}
\newcommand{\IvrVixMean}{17.1}
\newcommand{\IvrVixRvMean}{14.7}
\newcommand{\IvrVixVarRatioCovid}{0.73}
\newcommand{\IvrVixVarRatioEx}{1.36}
\newcommand{\OptDailyJ}{57.7}
\newcommand{\OptMeanPnl}{18}
\newcommand{\OptOptions}{1{,}898}
\newcommand{\OptOwnVsVix}{\ensuremath{-4.3}\%}
\newcommand{\OptOwnVsVixz}{\ensuremath{-1.5}}
\newcommand{\OptPremium}{179}
\newcommand{\ParAlpha}{0.527}
\newcommand{\ParHurst}{0.027}
\newcommand{\ParVolZero}{8.9}
\newcommand{\ParVzero}{0.0079}
\newcommand{\RegMSRatio}{1.74}
\newcommand{\RegMSRatioLo}{1.40}
\newcommand{\RegMSsigCalm}{0.71}
\newcommand{\RegMSsigStress}{1.23}
\newcommand{\RegRvQ}{17.4}
\newcommand{\RegVixQ}{20.16}
\newcommand{\RegVixQlo}{17.88}
\newcommand{\RepBS}{0.732}
\newcommand{\RepBSStored}{0.726}
\newcommand{\RepGKW}{0.828}
\newcommand{\RepGKWStored}{0.824}
\newcommand{\SimFlatSteps}{62}
\newcommand{\SimGKWvsBSzMax}{\ensuremath{-6.2}}
\newcommand{\SimGKWvsBSzMin}{\ensuremath{-5.1}}
\newcommand{\SimLowGainMax}{31\%}
\newcommand{\SimLowGainMin}{6\%}
\newcommand{\SimLowHoldoutAgree}{15/16}
\newcommand{\SimOptBSChange}{\ensuremath{-5.9}\%}
\newcommand{\SimOptBSz}{\ensuremath{-1.7}}
\newcommand{\SimOptGKWChange}{\ensuremath{-1.2}\%}
\newcommand{\SimOptGKWz}{\ensuremath{-0.5}}
\newcommand{\SimPaths}{50{,}000}
\newcommand{\SimSteps}{512}
\newcommand{\SynDailyJ}{64.0}
\newcommand{\SynMeanPnl}{30}
\newcommand{\SynOptions}{1{,}914}
\newcommand{\SynPremium}{197}
  % every number comes from src/report_assets.py
% -----------------------------------------------------------------------------

\begin{document}
\title{Viva Note: Variance-Optimal Hedging under Rough Volatility and Jumps}
\author{MAC F244 project team \\ Internal study note (not for submission)}
\date{\today}
\maketitle
\vspace{-1em}

\noindent This note explains, in plain language, what we did and why (main
reference: Shreve, \emph{Stochastic Calculus for Finance II}, Ch.~4--6 and 11).
Sections 2--4 are the mathematics to know; Section 5
describes the simulation model at the level needed; Section 6 gives our
results; Section 8 prepares for questions. Every number is generated from our
result files, so it matches the report.

\heading{1. The question in one paragraph.}
\noindent You sell a call option and want to remove its risk by trading the
underlying. In Black--Scholes this works perfectly if you rebalance
continuously and for free. Real hedgers trade at most a few times a day, pay
costs, and face markets that switch between calm and stress. We ask how often
to rebalance, whether to trade more in stress, and how much a better hedge
ratio helps. We test this in simulation, under a realistic model from Wang and
Wang's paper using their code, and on NIFTY data with synthetic and traded
options.

\heading{2. Black--Scholes delta hedging.}
\noindent Under $\mathbb{Q}$ with zero rate, $dS=\sigma S\,dW$. For the call price
$C(t,S)$, It\^o's formula gives
\[
dC=\big(C_t+\tfrac12\sigma^2S^2C_{SS}\big)\,dt+C_S\,dS .
\]
The Black--Scholes--Merton PDE says the bracket is zero, so $dC=C_S\,dS$:
holding $\Delta=C_S$ shares replicates the call. Equivalently, by the
martingale representation theorem, every $\mathbb{Q}$-martingale (here the
option price) is a constant plus a stochastic integral against the Brownian
motion, and hence against $S$. One source of risk and one traded asset: the market is complete.

\heading{3. What changes in practice (derivations to know).}
\noindent\emph{(a) Discrete rebalancing.} Hold $\Delta_k$ shares over an interval
of length $h$, during which the stock returns $R_k=\Delta S/S$. The short call's
P\&L is $\Delta_k\,\Delta S-\Delta C$. Expand
$\Delta C\approx C_t\,h+C_S\,\Delta S+\frac12\Gamma\,\Delta S^2$, with
$\Gamma=C_{SS}$, and use $C_t=-\frac12\sigma^2S^2\Gamma$ from the PDE:
\[
\varepsilon_k\approx\tfrac12\Gamma_kS_k^2\big(\sigma^2h-R_k^2\big).
\]
Since $\mathbb{E}[R_k^2]\approx\sigma^2h$, the mean is about zero. Since
$\mathrm{Var}(R_k^2)\approx2\sigma^4h^2$,
\[
\mathrm{Var}(\varepsilon_k)\approx\tfrac12\Gamma_k^2S_k^4\sigma^4h^2 .
\]
Errors in different intervals are uncorrelated (martingale differences), so
over the $T/h$ intervals the extra variance is of order $h$: halving the
interval roughly halves it. Our test suite checks this in a Black--Scholes
simulation.

\noindent\emph{(b) Gamma P\&L.} If the delta is computed at an implied
volatility $\sigma_i$ but the stock moves with realised volatility $\sigma_r$,
then $\mathbb{E}[R_k^2]\approx\sigma_r^2h$, and the expected P\&L per interval is
\[
\tfrac12\,\Gamma S^2\big(\sigma_i^2-\sigma_r^2\big)\,h .
\]
The seller gains when implied volatility exceeds the volatility that follows.
\emph{Our check (test period 2019--2026):}
\begin{compactitem}
\item India VIX was above the next month's realised volatility on \IvrVixAbove\%
of days (mean \IvrVixMean\% against \IvrVixRvMean\%).
\item Traded options' implied vol was above realised vol over their life for
\IvrOptAbove\% of options (mean \IvrOptIvMean\% against \IvrOptRvMean\%).
\item In variance terms (what the formula uses), implied exceeded realised by a
factor \IvrVixVarRatioEx{} (VIX) and \IvrOptVarRatioEx{} (options) outside the
COVID window. The crash reversed it (\IvrVixVarRatioCovid{} and
\IvrOptVarRatioCovid).
\end{compactitem}

\noindent\emph{(c) Costs.} Each trade costs $\kappa$ times the notional traded.
Trading more often reduces the error in (a) but raises costs: a trade-off. With
NIFTY's costs (up to 20\,bp), daily rebalancing still wins.

\heading{4. Incomplete markets and the minimum-variance hedge.}
\noindent\emph{Why the market becomes incomplete.} With stochastic volatility
(a second Brownian motion) or jumps of random size, there are more sources of
risk than traded assets, so the martingale representation in terms of $S$
alone fails and no stock position replicates the call.

\noindent\emph{It\^o's formula with jumps} (Shreve Ch.~11). For a process $X$
with jumps $\Delta X_s=X_s-X_{s-}$,
\begin{align*}
f(X_t)={}&f(X_0)+\int_0^t f'(X_{s-})\,dX_s+\frac12\int_0^t f''(X_{s-})\,d[X]^c_s\\
&+\sum_{0<s\le t}\big(f(X_s)-f(X_{s-})-f'(X_{s-})\Delta X_s\big).
\end{align*}
The sum replaces the linear (delta) approximation by the exact change of $f$ at
each jump. That gap is why a delta hedge cannot cover jumps: a jump moves the
option and the stock by different relative amounts.

\noindent\emph{The minimum-variance hedge.} Minimise
$\mathbb{E}^{\mathbb{Q}}[(Y-x-\int_0^T\vartheta_t\,dS_t)^2]$ for the payoff
$Y=(S_T-K)^+$. Let $M_t=\mathbb{E}^{\mathbb{Q}}[Y\mid\mathcal F_t]$ (the option price) and
decompose
\[
Y=\mathbb{E}^{\mathbb{Q}}[Y]+\int_0^T\vartheta^*_t\,dS_t+L_T,\qquad
\vartheta^*_t=\frac{d\langle M,S\rangle_t}{d\langle S\rangle_t},
\]
where the residual $L$ is uncorrelated with $S$ (the Galtchouk--Kunita--Watanabe
decomposition). Think of it as a regression: $\vartheta^*$ is the ``beta'' of
option moves on stock moves, and $L$ is what no stock position can explain. The
three pieces (a constant, a mean-zero stochastic integral and a residual
uncorrelated with $S$) are orthogonal, so by Pythagoras, for any
$(x,\vartheta)$,
\[
\mathbb{E}\Big[\Big(Y-x-\int\vartheta\,dS\Big)^2\Big]
=(\mathbb{E}[Y]-x)^2+\mathbb{E}\int_0^T(\vartheta^*-\vartheta)^2\,d\langle S\rangle+\mathbb{E}[L_T^2].
\]
So the best initial capital is the price, the best holding is $\vartheta^*$, and
$\mathbb{E}[L_T^2]$ is the unhedgeable risk.

\noindent\emph{Example to derive on the board (Heston).} Let $dS=\sqrt v\,S\,dW$
and $dv=\kappa_v(\theta-v)\,dt+\xi\sqrt v\,dB$ with $d\langle W,B\rangle=\rho\,dt$,
and $M_t=C(t,S_t,v_t)$. Two-variable It\^o gives
$dM=C_S\,dS+C_v\,dv+(\dots)\,dt$, so
$d\langle M,S\rangle=(C_S\,vS^2+C_v\,\rho\xi vS)\,dt$ and $d\langle S\rangle=vS^2\,dt$:
\[
\vartheta^*=C_S+\frac{\rho\,\xi\,C_v}{S}.
\]
With $\rho<0$ (volatility rises when prices fall) and $C_v>0$, you hold fewer
shares than the delta: a price fall comes with a volatility rise that partly
protects the option's value. In jump models a similar covariance term appears
for the jumps.

\heading{5. The simulation model.}
\noindent The paper's rough Hawkes--Heston model (under $\mathbb{Q}$, zero rate) is
\[
\frac{dS_t}{S_{t-}} = \sqrt{V_t}\,dW^S_t + (\text{compensated crash jumps}),
\]
\[
V_t = g_0(t) + \int_0^t K(t-s)\,dZ_s,\qquad K(t)=\frac{t^{\alpha-1}}{\Gamma(\alpha)},
\]
where $Z$ collects the variance's own shocks (correlated with the price through
$\rho$) and the jumps.
\begin{compactitem}
\item \emph{Leverage.} $\rho<0$: falling prices come with rising volatility.
\item \emph{Crash jumps.} A jump of size $z$ multiplies the price by
$e^{-\Lambda z}<1$ and pushes the variance up.
\item \emph{Clustering.} Jumps arrive at a rate proportional to $V$, and each jump
raises $V$, so jumps come in clusters.
\item \emph{Memory and roughness.} $V$ depends on all past shocks with power-law
weights $K$; with $\alpha<1$ the paths are rough.
\item \emph{Compensated} jumps have their expected effect subtracted, so the
price has no drift under $\mathbb{Q}$.
\end{compactitem}

\begin{table}[h]
\centering\small
\begin{tabular}{llp{8.2cm}}
\toprule
Parameter & Value & Meaning \\
\midrule
$\alpha$ & $0.527$ & roughness of the kernel $K(t)=t^{\alpha-1}/\Gamma(\alpha)$ \\
$\rho$ & $-0.731$ & correlation between price and variance shocks (leverage) \\
$b$ & $-1.812$ & mean reversion of the variance (negative = pulls back) \\
$c$ & $0.115$ & vol-of-vol: size of Brownian shocks to the variance \\
$\Lambda$ & $0.276$ & a jump of mark $z$ moves the log-price by $-\Lambda z$ \\
$\beta$ & $0.049$ & slope of the initial curve $g_0(t)=V_0+\beta t^\alpha/\Gamma(\alpha+1)$ \\
$V_0$ & $0.0079$ & initial variance \\
$\kappa$, $m$ & $1, 1$ & jump marks: $\nu(dz)=\kappa e^{-z/m}dz/m$ \\
$S_0$ & $1.0$ & initial price (prices are relative to $S_0$) \\
\bottomrule
\end{tabular}

\caption{Model parameters used in the paper and its code (initial volatility $\sqrt{V_0}\approx\ParVolZero\%$).}
\end{table}

\noindent\emph{How the hedge is computed.} Wang and Wang show that $\vartheta^*$
has the covariance form of Section 4 (a Brownian part involving $\rho$, plus a
jump part). They compute it with Fourier methods and a Riccati--Volterra
equation. We use their published code and verified that it reproduces their
results.

\heading{6. What we did and what we found.}
\noindent\emph{(a) Replication.} The official code gives a variance reduction
against no hedge of \RepGKWStored{} (GKW) and \RepBSStored{} (BS delta); our
rerun gives \RepGKW{} and \RepBS{}. The difference is Monte Carlo noise.
\emph{Careful:} these two numbers come from the official code, not from the
arXiv paper.

\begin{table}[h]
\centering\small
\begin{tabular}{lcc}
\toprule
Hedge & Official output & Our rerun \\
\midrule
Variance-optimal (GKW) hedge & 0.824 & 0.828 \\
Black--Scholes delta (implied vol) & 0.726 & 0.732 \\
\bottomrule
\end{tabular}

\caption{Replication: variance reduction vs.\ no hedge.}
\label{tab:rep}
\end{table}

\noindent\emph{(b) Discrete rebalancing.} If you trade only at
$t_0<\dots<t_{N-1}$ and pay $\kappa$ per unit of notional traded, the P\&L of
the hedged short call is
\[
e = x + \sum_k\varphi_k(S_{k+1}-S_k) - \kappa\sum_k S_k|\varphi_k-\varphi_{k-1}|
 - \kappa S_N|\varphi_{N-1}| - Y .
\]

\begin{table}[h]
\centering\small
\begin{tabular}{rcc}
\toprule
Rebalancing dates & GKW hedge & BS delta \\
\midrule
512 & 0.828 (0.012) & 0.732 (0.026) \\
64 & 0.806 (0.013) & 0.670 (0.031) \\
16 & 0.738 (0.016) & 0.541 (0.033) \\
4 & 0.583 (0.021) & 0.333 (0.036) \\
1 & 0.374 (0.019) & 0.108 (0.036) \\
\bottomrule
\end{tabular}

\caption{Simulation: variance reduction (s.e.) against the number of rebalancing dates.}
\end{table}

\noindent Findings.
\begin{compactitem}
\item Error grows as trading becomes rarer, as Section 3(a) predicts.
\item The GKW advantage over BS grows with rarer trading: the better hedge ratio
matters more when the position is held longer.
\item GKW beats BS at every cost level ($z$ from \SimGKWvsBSzMin{} to
\SimGKWvsBSzMax{}).
\item Smart rules (bands, variance-triggered) beat fixed schedules by
\SimLowGainMin--\SimLowGainMax{} at the same expected cost, but only when
turnover is low. At the cost-optimal point the change is only
\SimOptGKWChange{} for GKW ($z=\SimOptGKWz$), which is not significant.
\end{compactitem}

\noindent\emph{(c) Data and regimes.} The data are official NIFTY~50 closes and
India VIX checked against NSE. Regime labels (calm/stress) are calibrated on
2010--2018 only and use data up to day $t$ only. The main label is VIX above
\RegVixQ{} (with hysteresis down to \RegVixQlo); the second is 21-day realised
volatility above \RegRvQ\%. A statistical Markov-switching label is used only as a
robustness check. Out of sample, stress days are followed by roughly twice the
volatility (Table~\ref{tab:vreg}).

\begin{table}[h]
\centering\small
\begin{tabular}{lcc}
\toprule
Variant & VIX (hysteresis) & 21-day realised vol \\
\midrule
Main (train to 2018) & 1.94 [1.58, 2.44] & 1.79 [1.37, 2.41] \\
Train to 2017 & 1.94 [1.57, 2.46] & 1.80 [1.37, 2.38] \\
Train to 2019 & 2.01 [1.65, 2.54] & 1.85 [1.40, 2.47] \\
Start 2008 (unverified VIX) & 2.25 [1.52, 3.35] & 1.86 [1.32, 2.64] \\
\bottomrule
\end{tabular}

\caption{Test period: future 5-day volatility after stress / after calm [95\% CI].}
\label{tab:vreg}
\end{table}

\noindent\emph{(d) Backtests on NIFTY.} We sell a one-month ATM call each day and
delta-hedge it. Options are either synthetic (priced at VIX) or traded (NSE
bhavcopy, hedged with futures).
\begin{compactitem}
\item Daily rebalancing is best for costs up to 20\,bp, and the regime rule picked
in training is worse in testing.
\item At equal turnover, regime rules win only because of the COVID crash.
\item Inverted labels (trading less in stress) are clearly worse.
\item The seller's average gain is the gamma P\&L of Section 3(b).
\end{compactitem}

\begin{table}[h]
\centering\small
\begin{tabular}{lccccc}
\toprule
 & \multicolumn{2}{c}{Picked on train vs daily} & \multicolumn{3}{c}{Equal turnover, 5\,bp} \\
\cmidrule(lr){2-3}\cmidrule(lr){4-6}
Variant & 5\,bp & 20\,bp & beat fixed & med.\ $z$ & beat placebo \\
\midrule
Own implied vol & \ensuremath{+35.0}\% (\ensuremath{2.1}) & \ensuremath{+15.1}\% (\ensuremath{1.4}) & 48/48 & \ensuremath{-1.0} & 48/48 \\
Own IV, excl.\ COVID window & \ensuremath{+51.8}\% (\ensuremath{2.0}) & \ensuremath{+18.8}\% (\ensuremath{1.3}) & 17/48 & \ensuremath{0.4} & 41/48 \\
India VIX delta & \ensuremath{+33.7}\% (\ensuremath{2.1}) & \ensuremath{+16.3}\% (\ensuremath{1.5}) & 47/48 & \ensuremath{-1.0} & 46/48 \\
\bottomrule
\end{tabular}

\caption{Traded-option backtest (test period).}
\end{table}

\heading{7. How we made the results trustworthy.}
\begin{compactitem}
\item \emph{Same paths for all rules} (common random numbers) and \emph{paired}
standard errors. Jumps make unpaired errors huge.
\item \emph{Equal-turnover comparison.} A rule that trades more always has lower
error and higher cost, so we compare a rule with a coin-flip mixture of the two
neighbouring fixed schedules that has the same expected turnover. The mixture is
itself a valid strategy, so this is a fair benchmark.
\item \emph{No look-ahead.} Every decision at $t$ uses data up to $t$. This is
tested by truncating or scrambling future data and checking that past
decisions do not change.
\item \emph{Train/test split.} Rules and thresholds are chosen on 2010--2018 and
evaluated once on 2019--2026.
\item \emph{Placebos.} Inverted labels and year-old labels show whether any gain
comes from the regime information.
\item \emph{Block bootstrap.} Options written on consecutive days overlap, so
their errors are correlated. Resampling blocks of days keeps that dependence.
\item \emph{Criterion $J=\mathrm{Var}(e)+(\mathbb{E}\,\text{cost})^2$ on real
data.} The seller's mean gain (about $+\OptMeanPnl$\,bp per traded option) is
gamma P\&L, not hedging quality, so we remove it. $J$ equals the simulation
criterion when the premium is fair.
\item \emph{Data checks.} Expiring futures settle at the index close at all
\FoExpiries{} expiries, put--call parity holds within $\pm\FoParityBand$\,bp,
and the VIX data were reconciled against NSE.
\item \emph{Tests.} An automated test suite (\texttt{tests/}), including a
hand-computed option and agreement with the official code to 15 digits.
\end{compactitem}

\heading{8. Likely viva questions.}
\qa{Why is Black--Scholes hedging perfect, and what breaks it here?}{It\^o plus
the PDE gives $dC=C_S\,dS$, so the delta replicates the option. In practice we
trade discretely (Section 3a), pay costs, and face stochastic volatility and
jumps, which make the market incomplete (Section 4).}
\qa{Why does the error grow as rebalancing becomes rarer?}{Between trades the
option moves nonlinearly (gamma) while the hedge is frozen. Each interval adds
variance $\approx\frac12\Gamma^2S^4\sigma^4h^2$, so the total grows like $h$.}
\qa{Why did the option seller make money on NIFTY?}{Gamma P\&L:
$\frac12\Gamma S^2(\sigma_i^2-\sigma_r^2)$. Implied volatility was above the
volatility that followed most of the time (\IvrVixAbove\% of days for VIX), except
in the COVID crash.}
\qa{Why is the market incomplete?}{There are more sources of risk (a variance
shock and jumps of random size) than traded instruments (the stock).}
\qa{Why is the optimal initial capital the option price?}{The gain
$\int\vartheta\,dS$ has mean zero under $\mathbb{Q}$, so the error splits off
$(\mathbb{E}[Y]-x)^2$, which is minimised at $x=\mathbb{E}^{\mathbb{Q}}[Y]$.}
\qa{What is $\langle M,S\rangle$?}{The expected instantaneous covariance of the
option price and the stock. The holding must be chosen before the move, so it
uses this expected covariance, not the realised one.}
\qa{Did you derive the paper's hedge formula?}{No. We used the authors'
published implementation and checked that it reproduces their results. Our
contribution is the analysis of discrete, costly and regime-aware hedging, and
the NIFTY backtests.}
\qa{Why not just compare RMSE across rules?}{Rules with more trading have both
lower error and higher cost. Comparing at equal expected turnover (equal cost)
separates timing skill from simply trading more.}
\qa{Why the block bootstrap?}{Daily-written one-month options overlap by up to
20 days, so their errors are strongly correlated; i.i.d.\ standard errors would
be far too small.}
\qa{What does the COVID result mean?}{The apparent advantage of regime-aware
rebalancing comes from one extreme episode. Excluding it, the advantage
disappears. One crisis is too little evidence for a general effect.}
\qa{Why not fit the paper's model to NIFTY?}{That means fitting the whole
volatility smile with a specialised pricer: a project in itself. On real data we
therefore use Black--Scholes and Black--76 deltas.}
\qa{What would you do next?}{Make costs depend on the regime, compare the
minimum-variance hedge of Section 4 with the delta on real options, and add
static option hedges.}

\heading{9. Glossary.}
\noindent\emph{Delta, gamma}: first and second derivative of the option price
in the stock price.
\emph{Implied / realised volatility}: the volatility in the option's price /
the volatility the stock actually showed.
\emph{Variance reduction}: $1-\mathrm{Var}(\text{hedged})/\mathrm{Var}(\text{unhedged})$.
\emph{Turnover}: notional traded. \emph{bp}: basis point, $0.01\%$.
\emph{Placebo}: the same rule fed meaningless information.
\emph{Look-ahead bias}: using data not yet available at decision time.
\emph{Common random numbers}: evaluating all strategies on identical simulated paths.
\emph{Bhavcopy}: NSE's daily end-of-day file of all derivatives contracts.

\end{document}
